% ==============================================================
\section{実験}\label{sec:experiments}
% ==============================================================

以下の研究課題を検証する。
\textbf{RQ1}: 1次元設定（空間次元なし）と比べて3次元設定は空間的偽陽性を抑制できるか？
\textbf{RQ2}: データサイズ $N$ に対して計算時間はどうスケールするか？
\textbf{RQ3}: 各高速化技術はどの程度の効果をもたらすか？（アブレーション）
\textbf{RQ4}: 既存手法（FIM、1次元 AW、Span\_3D、MZoom\_3D）と比べてどの程度優れているか？
\textbf{RQ5}: 次元数 $d$ の増加に対して計算時間はどうスケールするか？

\subsection{実験環境}

すべての実験は Rust（release build、Rayon 並列化）で実装した
\texttt{md\_pipeline} バイナリを用い、
Apple M3（8 コア、24 GB RAM）上で実行した。
デフォルトパラメータ：$\theta = 3$（サポート閾値）、
$W_t = 5$（時間ウィンドウ）、$W_x = W_y = 3$（空間ウィンドウ）、
$L_{\max} = 3$（最大パターン長）、$\Delta t = \Delta x = \Delta y = 1.0$。

\subsection{EX-A: 3次元 vs 1次元相当の比較}
\label{sec:ex_a}

\paragraph{目的.}
Apriori-Window~\cite{aprioriwindow}は時間軸のみを扱うため、
空間的に局在するパターンを識別できない（全空間を集約する）。
本実験では合成データを用いて、
3次元設定が空間的偽陽性を抑制することを実証する。

\paragraph{合成データ設計.}
$N = 5{,}000$ トランザクション、アイテム数 $|{\mathcal{I}}| = 10$、
空間範囲 $[0, 50]^2$、時刻範囲 $[0, 100]$ の合成データを生成する。
3つの密集領域を植え付ける：各領域は2アイテムパターン
$P_1 = \{1,2\}$、$P_2 = \{3,4\}$、$P_3 = \{5,6\}$ を割り当て、
それぞれランダムな時空間中心（$t_c, x_c, y_c$）の近傍に
20件の集中トランザクションを注入する。
残りの $N - 60$ 件はバックグラウンドノイズ（ランダム）。

\paragraph{比較設定.}
\begin{itemize}
\item \textbf{3次元設定}（本提案）：$W_x = W_y = 3$（植え付け領域に対応する局所ウィンドウ）
\item \textbf{1次元相当設定}：$W_x = W_y = 50$（空間範囲全体をカバーする大ウィンドウで
  空間次元を事実上無視し、1次元集約と同等の挙動を実現）
\end{itemize}
どちらも $\theta = 3$、$W_t = 5$ を使用する。

\paragraph{評価指標.}
検出された $k=2$ パターンの集合と植え付けパターン集合の一致で
Precision・Recall を計算する。5シードで実験する。

\paragraph{結果.}
表~\ref{tab:ex_a}に結果を示す。

\begin{table}[t]
\centering
\caption{EX-A: 3次元設定 vs 1次元相当設定の比較（5シード、$k=2$ パターン）。}
\label{tab:ex_a}
\small
\begin{tabular}{lrrcc}
\toprule
設定 & 検出 $k{\geq}2$ パターン数 & 植え付け回収数 & Precision & Recall \\
\midrule
3次元（$W_x{=}W_y{=}3$）     & 3   & 3 & 1.00 & 1.00 \\
1次元相当（$W_x{=}W_y{=}50$）& $\approx$93--100 & 3 & $\approx$0.03 & 1.00 \\
\bottomrule
\end{tabular}
\end{table}

\paragraph{分析.}
(RQ1)
3次元設定では全5シードで植え付け3パターン $\{1,2\}$、$\{3,4\}$、$\{5,6\}$ を
完全回収し、偽陽性パターンは0件であった（Precision $= 1.0$、Recall $= 1.0$）。
一方、1次元相当設定では $\{1,3\}$、$\{1,4\}$、$\{2,3\}$ などの
バックグラウンドノイズ由来の偽陽性パターンが多数（93〜100件）検出された
（Precision $\approx 0.03$）。

これは、空間次元を集約すると空間的に異なる位置のトランザクションが
同一ウィンドウに混入し、偶発的な共起が閾値を超えることを示す。
空間次元の明示的なモデル化が精度向上に本質的に重要である。

\subsection{EX-B: データサイズに対するスケーラビリティ}
\label{sec:ex_b}

\paragraph{実験設計.}
$N \in \{1{,}000,\; 5{,}000,\; 10{,}000,\; 50{,}000,\; 100{,}000\}$ と
ウィンドウサイズ $W_x = W_y \in \{1, 2, 3, 5, 7, 10\}$（$N = 20{,}000$ 固定）を変化させ、
実行時間を計測する。

\paragraph{結果.}
表~\ref{tab:ex_b_n}（$N$ 方向）と表~\ref{tab:ex_b_w}（ウィンドウ方向）に結果を示す。

\begin{table}[t]
\centering
\caption{EX-B: $N$ に対するスケーラビリティ（$W_x{=}W_y{=}3$、$W_t{=}5$）。}
\label{tab:ex_b_n}
\small
\begin{tabular}{rrc}
\toprule
$N$ & 検出パターン数 & 実行時間 (s) \\
\midrule
1,000   & 9  & 0.005 \\
5,000   & 13 & 0.021 \\
10,000  & 13 & 0.030 \\
50,000  & 55 & 0.455 \\
100,000 & 59 & 1.407 \\
\bottomrule
\end{tabular}
\end{table}

\begin{table}[t]
\centering
\caption{EX-B: 空間ウィンドウサイズに対するスケーラビリティ（$N{=}20{,}000$、$W_t{=}5$）。}
\label{tab:ex_b_w}
\small
\begin{tabular}{rrc}
\toprule
$W_x = W_y$ & 検出パターン数 & 実行時間 (s) \\
\midrule
1  & 13 & 0.010 \\
2  & 13 & 0.034 \\
3  & 22 & 0.086 \\
5  & 55 & 0.477 \\
7  & 55 & 1.150 \\
10 & 86 & 1.991 \\
\bottomrule
\end{tabular}
\end{table}

\paragraph{分析.}
(RQ2) 表~\ref{tab:ex_b_n}より、実行時間は $N$ に対してほぼ線形にスケールする。
$N = 100{,}000$ でも約 1.5 秒を達成しており、大規模データでの実用性を示す。

表~\ref{tab:ex_b_w}より、空間ウィンドウサイズの増大に対しては
超線形のスケーリングが観察される。
$W_x = W_y = 1$ から $W_x = W_y = 3$ まではほぼ $W^2$ に比例するが（0.010s → 0.086s、$8.6\times$）、
$W_x = W_y = 7$ 以上では XY イベントリスト構築の $O(N W_x W_y)$ コストが支配的となる。
この結果は計算量分析（表~\ref{tab:complexity}）と整合的である。

\subsection{EX-C: 高速化技術のアブレーション}\label{sec:ex_c}

\paragraph{実験設計.}
Apriori 枝刈りの効果を測定するため、
$L_{\max}$ を 1 から 2、3、4 と増やしながら実行時間を計測する。
真の頻出パターンが $k=2$ で打ち止まる合成データ（$N = 20{,}000$）を使用する。

\paragraph{結果と分析.}
(RQ3) $L_{\max}$ を 1 から 2、3、4 と増やしても
実行時間はほぼ不変（0.062→0.083→0.082→0.087 s）であった。
これは Apriori 枝刈りが $k \geq 3$ の候補を完全に排除していることを示す。
当該データでは $k = 3$ を満たすパターンが存在せず、
$\mathcal{C}(P) = \emptyset$ の枝刈りにより $I_P$ 構築が省略されたためである。
この挙動は Apriori-Window~\cite{aprioriwindow}の1次元枝刈り効果と同様であり、
3次元への拡張でも同じ恩恵が得られることを示す。

\subsection{EX-D: 比較手法との性能比較}
\label{sec:ex_d}

\paragraph{目的.}
本実験では5種類の手法を合成データ上で比較し、
3次元 Apriori-Window の有効性を検証する（RQ4）。

\paragraph{比較手法.}
\begin{itemize}
\item \textbf{FIM\_global}: グローバル頻出アイテムセットマイニング~\cite{agrawal1994apriori}。
  空間・時間構造を無視し、全トランザクション上の共起頻度のみを使用する。
\item \textbf{1D\_AW}: Apriori-Window~\cite{aprioriwindow}（1次元、時間軸のみ）。
  空間次元を集約し、時刻スライディングウィンドウでサポートを計算する。
\item \textbf{Span\_3D}: 空間的点密度手法。
  $(1 \times 1 \times 1)$ セルごとの共起頻度が閾値を超えるパターンを検出する。
\item \textbf{MZoom\_3D}: 密部分テンソル検出~\cite{shin2018mzoom}の3次元近似。
  時空間グリッドブロック内の最大密度に基づいてパターンを検出する。
\item \textbf{3次元 AW}（本提案）: 3次元スライディングウィンドウサポートを用いた提案手法。
\end{itemize}

\paragraph{実験設計.}
$N = 5{,}000$ トランザクション（アイテム数 $|\mathcal{I}| = 20$、空間範囲 $[0, 50]^2$、
時刻範囲 $[0, 100]$）の合成データを生成し、
3つの密集領域を植え付ける（各領域 $P_1 = \{1,2\}$、$P_2 = \{3,4\}$、$P_3 = \{5,6\}$、
各20件注入、ジッタ：$t \pm 2.0$、$x \pm 1.0$、$y \pm 1.0$）。
バックグラウンドノイズは全アイテム（1--20）のランダム組み合わせ。

全手法でグリッドサーチを実施し（パラメータ範囲を表~\ref{tab:ex_d_grid}に示す）、
F1 を最大化するパラメータを選択する。
5シードで実験し平均を報告する。

\begin{table}[t]
\centering
\caption{EX-D: グリッドサーチのパラメータ範囲。}
\label{tab:ex_d_grid}
\small
\begin{tabular}{lp{6cm}}
\toprule
手法 & 探索パラメータ \\
\midrule
FIM\_global   & $\theta \in \{3,5,7,10,15,20,30\}$ \\
1D\_AW        & $W_t \in \{3,5,7,10\}$、$\theta \in \{3,5,7,10\}$ \\
Span\_3D      & $\theta \in \{3,5,7,10,15\}$ \\
MZoom\_3D     & $\theta \in \{3,5,7\}$、最小密度 $\in \{0.1,0.3,0.5\}$ \\
3次元 AW      & $W_t \in \{3,5,7,10\}$、$W_x{=}W_y \in \{2,3,5,7\}$、$\theta \in \{3,5,7,10\}$ \\
\bottomrule
\end{tabular}
\end{table}

\paragraph{結果.}
表~\ref{tab:ex_d}に各手法の最良パラメータにおける検出性能を示す。

\begin{table}[t]
\centering
\caption{EX-D: 比較実験結果（合成データ、5シード平均±標準偏差）。
植え付け $k=2$ パターン3件の回収率で評価。}
\label{tab:ex_d}
\small
\begin{tabular}{lcccc}
\toprule
手法 & Precision & Recall & F1 & 検出パターン数 \\
\midrule
FIM\_global  & $0.016 \pm 0.000$ & $\mathbf{1.000} \pm 0.000$ & $0.031 \pm 0.000$ & 190.0 \\
1D\_AW       & $0.950 \pm 0.100$ & $\mathbf{1.000} \pm 0.000$ & $0.971 \pm 0.057$ & 3.2   \\
Span\_3D     & $\mathbf{1.000} \pm 0.000$ & $0.733 \pm 0.133$ & $0.840 \pm 0.080$ & 2.2   \\
MZoom\_3D    & $\mathbf{1.000} \pm 0.000$ & $\mathbf{1.000} \pm 0.000$ & $\mathbf{1.000} \pm 0.000$ & 3.0 \\
\textbf{3次元 AW}（本提案） & $\mathbf{1.000} \pm 0.000$ & $\mathbf{1.000} \pm 0.000$ & $\mathbf{1.000} \pm 0.000$ & \textbf{3.0} \\
\bottomrule
\end{tabular}
\end{table}

\paragraph{分析.}
(RQ4)
\textbf{FIM\_global} は全3パターンを回収（Recall $= 1.0$）するが、
190件の偽陽性パターンを生成する（Precision $= 0.016$）。
これは空間・時間構造を無視することで、
バックグラウンドノイズ由来の偶発的共起が全てのパターンとして報告されるためである。

\textbf{1D\_AW}（Apriori-Window~\cite{aprioriwindow}）は時間軸のみを使用し、
大半のシードで完全回収を達成する（Recall $= 1.0$、F1平均 $= 0.971$）が、
1シード（seed=456）で偽陽性パターン1件が発生した（検出数4件）。
これは時刻が近接する異なる空間位置のトランザクションが同一時間ウィンドウに混入するためである。

\textbf{Span\_3D} は高精度（Precision $= 1.0$）だが、
平均 Recall $= 0.733$（3シードで2件回収、2シードで3件回収）。
ウィンドウ密度ではなく点密度を使用するため、
注入領域（ジッタ幅 $\pm 2$ の空間的広がり）の共起が分散した場合に閾値未達となる。

\textbf{MZoom\_3D}（近似実装）はブロック密度によるパターン検出であり、
本実験では F1 $= 1.0$ を達成した。
ただし MZoom の本来の目的はテンソル異常検出であり、
アイテムセットのスライディングウィンドウサポート意味論は持たない。

\textbf{3次元 AW}（本提案）は全シードで Precision $= 1.0$、Recall $= 1.0$ を達成し、
唯一の手法として常に正確な3パターン検出を実現した。
時空間スライディングウィンドウサポートにより、
空間的に局在した密集アイテムセット領域を偽陽性なく発見できる。

\subsection{EX-E: 次元数スケーラビリティ}\label{sec:ex_e}

\paragraph{目的.}
次元数 $d$ の増加が計算時間に与える影響を定量化する（RQ5）。

\paragraph{実験設計.}
$N = 10{,}000$ トランザクション固定、
$d \in \{2, 3, 4, 5\}$ を変化させる。
各次元のビンサイズ $\Delta_k = 1.0$、
ウィンドウサイズ $W_k = 3$（全次元統一）、
閾値 $\theta = 3$、$L_{\max} = 2$。
座標範囲は各次元 $[0, 50]$ に統一し、
密集領域は3件植え付ける（各次元の中心・ジッタを適宜設定）。
5シードの平均実行時間を計測する。

\paragraph{結果.}
表~\ref{tab:ex_e}に結果を示す。

\begin{table}[t]
\centering
\caption{EX-E: 次元数 $d$ に対するスケーラビリティ
（$N{=}10{,}000$、$W_k{=}3$、$\theta{=}3$、5シード平均）。}
\label{tab:ex_e}
\small
\begin{tabular}{rrrr}
\toprule
$d$ & グリッドサイズ $G$ & 実行時間 (s) & 備考 \\
\midrule
2 & $50 \times 50$           & \textit{TBD} & 1D-AW 相当 \\
3 & $50 \times 50 \times 50$ & \textit{TBD} & 本稿 EX-B 相当 \\
4 & $50^4$                   & \textit{TBD} & \\
5 & $50^5$                   & \textit{TBD} & \\
\bottomrule
\end{tabular}
\end{table}

\paragraph{分析.}
(RQ5)
\textit{[実装完了後に実測値で更新]}

\subsection{考察}\label{sec:discussion}

\paragraph{密集領域の数と空間ウィンドウ.}
$W_x = W_y$ の増大に伴い検出パターン数が増加する（表~\ref{tab:ex_b_w}）。
EX-A の結果と合わせると、空間ウィンドウを大きくするほど
偽陽性パターンが増加することがわかる。
適切な空間ウィンドウサイズはデータの空間粒度に依存し、
応用ドメインの知識に基づいた設定が重要である。
